\ifthenelse{\equal{\LanguageOption}{spanish}}{%
    \chapter*{Motivación}
    
    La motivación principal para desarrollar esta investigación surge de la necesidad urgente de abordar la vulnerabilidad de los menores de edad frente a las crecientes amenazas en el entorno digital. En los últimos años, la adopción masiva de dispositivos móviles y la inmersión en plataformas interactivas han expuesto a los niños a riesgos silenciosos pero devastadores, como el ciberacoso, la exposición a contenido explícito y el \textit{grooming}, afectando profundamente su bienestar psicológico y su seguridad.

    Como estudiante comprometido con el desarrollo de soluciones tecnológicas con impacto social, sentí la necesidad de aportar una alternativa moderna y ética frente a las herramientas de control parental tradicionales. Esto me impulsó a explorar e integrar tecnologías de vanguardia, como la Inteligencia Artificial de Borde (\textit{Edge AI}), el Procesamiento de Lenguaje Natural y la Visión Artificial. Esta problemática representó no solo un fascinante desafío técnico de optimización algorítmica en dispositivos móviles, sino también la oportunidad de demostrar que la tecnología puede proteger activamente sin vulnerar el derecho constitucional a la privacidad.

    Además, este proyecto responde a mi interés personal por aplicar los conocimientos en desarrollo de \textit{software} y \textit{Machine Learning} adquiridos durante mi formación académica en un estudio que tenga un impacto directo y positivo en las familias. La seguridad digital es un componente clave para el desarrollo integral en la sociedad contemporánea, y mejorarla significa salvaguardar la integridad de quienes son más vulnerables.

    En resumen, la presente investigación nace del compromiso de aportar una perspectiva técnica y fundamentada que redefina la ciberseguridad infantil, ofreciendo un sistema inteligente y autónomo que promueva un ecosistema digital más seguro, privado y confiable para las nuevas generaciones.
}{%
    \chapter*{Motivation}
    
    The primary motivation for developing this research arises from the urgent need to address the vulnerability of minors to growing threats in the digital environment. In recent years, the massive adoption of mobile devices and immersion in interactive platforms have exposed children to silent but devastating risks, such as cyberbullying, exposure to explicit content, and grooming, deeply affecting their psychological well-being and safety.

    As a student committed to developing technological solutions with social impact, I felt the need to provide a modern and ethical alternative to traditional parental control tools. This drove me to explore and integrate cutting-edge technologies, such as Edge Artificial Intelligence (Edge AI), Natural Language Processing, and Computer Vision. This problem represented not only a fascinating technical challenge of algorithmic optimization on mobile devices, but also the opportunity to demonstrate that technology can actively protect without violating the constitutional right to privacy.

    Furthermore, this project responds to my personal interest in applying the software development and Machine Learning knowledge acquired during my academic training into a study that has a direct and positive impact on families. Digital safety is a key component for holistic development in contemporary society, and improving it means safeguarding the integrity of the most vulnerable.

    In summary, this research is born from the commitment to provide a technical and well-founded perspective that redefines child cybersecurity, offering an intelligent and autonomous system that promotes a safer, more private, and reliable digital ecosystem for new generations.
}

\MediaOptionLogicBlank